\documentclass[11pt]{article}
\usepackage[margin=1in]{geometry}
\usepackage{amsmath,amssymb,amsthm,mathtools}
\usepackage{hyperref}
\usepackage{enumitem}
\usepackage{xcolor}
\hypersetup{colorlinks=true, linkcolor=blue, urlcolor=blue, citecolor=blue}

\newtheorem{theorem}{Theorem}
\newtheorem{proposition}{Proposition}
\newtheorem{lemma}{Lemma}
\newtheorem{corollary}{Corollary}
\newtheorem{assumption}{Assumption}
\newtheorem{remark}{Remark}
\newtheorem{example}{Example}
\newtheorem{definition}{Definition}

\newcommand{\Q}{\mathcal Q}
\newcommand{\A}{\mathcal A}
\newcommand{\D}{\mathcal D}
\newcommand{\DeltaK}{\Delta_K}
\newcommand{\E}{\mathbb E}
\newcommand{\Pp}{\mathbb P}
\newcommand{\one}{\mathbf 1}
\newcommand{\what}{\widehat w}
\newcommand{\wstar}{w^\star}
\newcommand{\wfix}{w}
\newcommand{\eps}{\varepsilon}
\newcommand{\norm}[1]{\left\lVert #1\right\rVert}
\newcommand{\TV}{\mathrm{TV}}
\newcommand{\KL}{\mathrm{KL}}
\newcommand{\Ber}{\mathrm{Ber}}
\newcommand{\class}{\mathcal P}
\newcommand{\Mhat}{\widehat M}
\newcommand{\nmax}{n_{\max}}
\newcommand{\mmin}{m_{\min}}

\title{Best-of-Infinity Ensembles:\\ Buffered Weight Learning and Adaptive Test-Time Allocation}
\author{Union theory note from buffered ERM and allocation notes}
\date{\today}

\begin{document}
\maketitle

\begin{abstract}
This union note combines two complementary theory notes for Best-of-Infinity style test-time ensembling.  Part~I studies the statistical problem of learning ensemble weights from $Q$ labeled benchmark questions when each model's answer distribution is estimated from $N$ generations; its main result is a buffered empirical-risk theorem with a two-sample-size rate.  Part~II studies the allocation problem after a weight vector is fixed: given a generation budget, how should test-time samples be spent across problems with different margins?  The two pieces should be read as a single paper spine: first learn or validate weights under double randomness, then allocate additional test-time compute according to margin difficulty.  The fixed-writeup system should preserve this hierarchy, avoid claiming unrun experiments, and treat the empirical program at the end as planned unless concrete result artifacts are supplied.
\end{abstract}

\section*{How this union note should drive the paper}
\begin{enumerate}[label=(\roman*)]
\item The main paper should present a single problem: Best-of-Infinity ensembling has two statistical bottlenecks, labeled questions $Q$ and per-question generations.  Weight learning and compute allocation are two stages of the same pipeline.
\item Part~I is the core learning theorem: buffered ERM gives a robust $Q,N$ decomposition and explains why margin alone does not imply fast $Q$ rates.
\item Part~II is the compute-allocation theorem: at a fixed weight vector, margins determine the price of resolving each problem, and adaptive allocation can improve the generation exponent relative to uniform allocation.
\item Experiments should be written only when backed by the code/results inventory.  Otherwise, the empirical program should be described as a concrete plan over saved generations and bootstrap simulations.
\end{enumerate}

\part{Buffered Weight Learning under Double Randomness}
\section{Purpose of this note}

The goal is to isolate the statistical structure in Best-of-Infinity style LLM ensembling.  There are two independent sample sizes:
\[
\begin{aligned}
Q &= \text{number of labeled benchmark questions},\\
N &= \text{number of generations per model per question}.
\end{aligned}
\]
The first controls how well we learn ensemble weights across problems.  The second controls how well we estimate each model's answer distribution on each problem.  This note gives a clean theorem for a margin-buffered empirical risk minimizer and then states an optional fast-rate theorem under an additional geometric identifiability condition.

The main message is:
\[
\boxed{
R(\what)-R(\wstar)
\lesssim
\sqrt{\frac{K\log(AQ)+\log(1/\delta)}{Q}}
+
\left(\frac{\log(KAQ/\delta)}{N}\right)^{\beta/2}.}
\]
The $Q$-term is ordinary statistical generalization over labeled questions.  The $N$-term is the finite-generation floor.  The latter can be faster than $N^{-1/2}$ because the objective is a thresholded sign/correctness objective, not estimation of a smooth scalar parameter.

\section{Setup and notation}

Let $q\in\Q$ denote a problem, drawn from a distribution $\D$.  Each problem has a gold answer $a^\star(q)$ and a finite candidate answer set $\A(q)$ with
\[
|\A(q)|\le A.
\]
There are $K$ language models.  Model $i$ induces a distribution over final extracted answers,
\[
p_i(a\mid q),\qquad a\in\A(q),\quad i=1,\dots,K.
\]
For a fixed problem $q$, a fixed model $i$, and a generation budget $N$, we observe i.i.d. final answers
\[
Y_{i,1}(q),\dots,Y_{i,N}(q)\sim p_i(\cdot\mid q),
\]
and form the empirical answer distribution
\[
\widehat p_i(a\mid q)
:=
\frac1N\sum_{n=1}^N \one\{Y_{i,n}(q)=a\}.
\]
It is convenient to write the $K$-vector of model probabilities as
\[
p_\cdot(a\mid q)
:=
\big(p_1(a\mid q),\dots,p_K(a\mid q)\big)\in[0,1]^K,
\]
and similarly for $\widehat p_\cdot(a\mid q)$.

An ensemble is a weight vector
\[
w=(w_1,\dots,w_K)\in\Delta_K
:=
\left\{w\in\mathbb R_+^K:\sum_{i=1}^K w_i=1\right\}.
\]
Its true score for answer $a$ on problem $q$ is
\[
s_q(a;w):=\langle w,p_\cdot(a\mid q)\rangle,
\]
and its empirical score is
\[
\widehat s_q(a;w):=\langle w,\widehat p_\cdot(a\mid q)\rangle.
\]

For every wrong answer $a\ne a^\star(q)$ define the answer-gap vector
\[
\Delta_{q,a}
:=
p_\cdot(a^\star(q)\mid q)-p_\cdot(a\mid q)\in[-1,1]^K,
\]
and its empirical version
\[
\widehat\Delta_{q,a}
:=
\widehat p_\cdot(a^\star(q)\mid q)-\widehat p_\cdot(a\mid q).
\]
The true margin of $w$ on $q$ is
\[
M_q(w)
:=
\min_{a\in\A(q),\,a\ne a^\star(q)}
\langle w,\Delta_{q,a}\rangle.
\]
The empirical margin is
\[
\widehat M_q(w)
:=
\min_{a\in\A(q),\,a\ne a^\star(q)}
\langle w,\widehat\Delta_{q,a}\rangle.
\]
Thus $M_q(w)>0$ means that the true weighted ensemble assigns higher probability to the gold answer than to every wrong answer.  We count ties as errors.

The clean loss and clean risk are
\[
\ell_q(w):=\one\{M_q(w)\le 0\},
\qquad
R(w):=\Pp_{q\sim\D}\{M_q(w)\le 0\}.
\]
Let
\[
\wstar\in\arg\min_{w\in\Delta_K} R(w)
\]
be a clean-risk minimizer.

The training data are
\[
\mathcal S_{Q,N}
=
\left\{\left(q_j,a^\star(q_j),\{Y_{j,i,n}:i\in[K],n\in[N]\}\right):j=1,\dots,Q\right\},
\]
where $q_j\overset{i.i.d.}{\sim}\D$.  On the training set, $\widehat M_{q_j}(w)$ is observable because the benchmark supplies $a^\star(q_j)$ and the generations supply $\widehat p_i(\cdot\mid q_j)$.

\section{The margin-buffered estimator}

Raw ERM minimizes
\[
\frac1Q\sum_{j=1}^Q \one\{\widehat M_{q_j}(w)\le 0\}.
\]
This can exploit finite-$N$ noise near the decision boundary.  We instead use a buffer:
\[
\widehat R_\eps(w)
:=
\frac1Q\sum_{j=1}^Q \one\{\widehat M_{q_j}(w)\le \eps\},
\]
and define
\[
\what_\eps\in\arg\min_{w\in\Delta_K}\widehat R_\eps(w).
\]
The buffer is chosen at the scale of the uniform finite-generation error on the training sample:
\[
\eps_N(\delta)
:=
\sqrt{\frac{2\log(4KAQ/\delta)}{N}}.
\]
This choice is not sacred; any constant multiple of the same order works.

\section{A basic concentration event}

\begin{lemma}[Uniform training-sample margin accuracy]
\label{lem:margin-accuracy}
Fix $\delta\in(0,1)$.  With probability at least $1-\delta/2$ over the generation randomness conditional on $q_1,\dots,q_Q$,
\[
\sup_{j\le Q}\sup_{w\in\Delta_K}
\left|\widehat M_{q_j}(w)-M_{q_j}(w)\right|
\le
\eps_N(\delta).
\]
\end{lemma}

\begin{proof}
It suffices to control all answer frequencies.  By Hoeffding's inequality and a union bound,
\[
\Pp\left(
\max_{j\le Q}\max_{i\le K}\max_{a\in\A(q_j)}
\left|\widehat p_i(a\mid q_j)-p_i(a\mid q_j)\right|
>
\frac{\eps_N}{2}
\right)
\le
2KAQ\exp\left(-\frac{N\eps_N^2}{2}\right)
\le \frac\delta2.
\]
On this event, for every wrong answer $a$,
\[
\left|\langle w,\widehat\Delta_{q_j,a}-\Delta_{q_j,a}\rangle\right|
\le
\eps_N
\]
because $w\in\Delta_K$.  Taking minima over $a$ preserves the same Lipschitz bound.
\end{proof}

On the event in Lemma~\ref{lem:margin-accuracy}, for every training problem and every $w$,
\[
\one\{M_{q_j}(w)\le 0\}
\le
\one\{\widehat M_{q_j}(w)\le \eps_N\}
\le
\one\{M_{q_j}(w)\le 2\eps_N\}.
\]
This is the key sandwich.  The empirical noisy-buffered loss lies between the clean zero-buffer loss and the clean $2\eps_N$-buffer loss.

\section{Capacity of the clean margin classes}

For $t\ge0$ define the thresholded clean-margin class
\[
\mathcal F_t
:=
\left\{q\mapsto \one\{M_q(w)\le t\}:w\in\Delta_K\right\}.
\]
For fixed problems $q_1,
\dots,q_Q$, each pair $(q_j,a)$ with $a\ne a^\star(q_j)$ contributes the affine hyperplane
\[
\langle w,\Delta_{q_j,a}\rangle=t
\]
in the $(K-1)$-dimensional simplex.  There are at most $Q(A-1)$ such hyperplanes.  Therefore the number of possible sign patterns over $q_1,\dots,q_Q$ is at most the number of cells in an arrangement of $Q(A-1)$ hyperplanes in dimension $K-1$:
\[
\Pi_{\mathcal F_t}(Q)
\le
\sum_{r=0}^{K-1}\binom{Q(A-1)}{r}
\le
\left(\frac{eQ(A-1)}{K-1}\right)^{K-1}
\]
when $Q(A-1)\ge K-1$.  Thus the effective combinatorial dimension is of order $K\log(AK)$, or equivalently the finite-sample deviation carries a factor $(K-1)\log(AQ)$, not $KA$.

We will use the shorthand
\[
\mathfrak v_Q(\delta)
:=
C
\sqrt{
\frac{(K-1)\log(eAQ)+\log(1/\delta)}{Q}
},
\]
where $C>0$ is a universal constant.  Standard VC symmetrization and the above growth bound imply that, with probability at least $1-\delta/2$ over $q_1,\dots,q_Q$,
\[
\sup_{w\in\Delta_K,\,t\in\{0,2\eps_N\}}
\left|
\frac1Q\sum_{j=1}^Q\one\{M_{q_j}(w)\le t\}
-
\Pp_q\{M_q(w)\le t\}
\right|
\le
\mathfrak v_Q(\delta).
\]

\section{Main theorem: buffered ERM}

\begin{assumption}[Comparator margin condition]
\label{ass:comp-margin}
There are constants $C_m>0$, $\beta>0$, and $t_0>0$ such that
\[
\Pp_{q\sim\D}\{0<M_q(\wstar)\le t\}
\le
C_m t^\beta,
\qquad 0<t\le t_0.
\]
\end{assumption}

This is a one-sided condition: it only controls correctly solved problems at $\wstar$ whose positive margin is small.  That is enough for the buffered theorem because the proof compares the learned rule to $\wstar$ and only pays for the extra $2\eps_N$ band around correct decisions of $\wstar$.

\begin{theorem}[Buffered ERM, slow $Q$-rate and fast finite-generation floor]
\label{thm:buffered-main}
Assume Assumption~\ref{ass:comp-margin}.  Let
\[
\eps_N=\sqrt{\frac{2\log(8KAQ/\delta)}{N}}
\]
and suppose $2\eps_N\le t_0$.  Let
\[
\what_\eps\in\arg\min_{w\in\Delta_K}
\frac1Q\sum_{j=1}^Q\one\{\widehat M_{q_j}(w)\le\eps_N\}.
\]
Then with probability at least $1-\delta$ over both the $Q$ problem draws and the $N$ generations per model per problem,
\[
R(\what_\eps)-R(\wstar)
\le
2\mathfrak v_Q(\delta)
+
C_m(2\eps_N)^\beta.
\]
Equivalently, up to universal constants,
\[
\boxed{
R(\what_\eps)-R(\wstar)
\lesssim
\sqrt{\frac{(K-1)\log(eAQ)+\log(1/\delta)}{Q}}
+
C_m
\left(\frac{\log(KAQ/\delta)}{N}\right)^{\beta/2}.
}
\]
\end{theorem}

\begin{proof}
Let
\[
\widetilde R_t(w):=\frac1Q\sum_{j=1}^Q\one\{M_{q_j}(w)\le t\},
\qquad
R_t(w):=\Pp_q\{M_q(w)\le t\}.
\]
On the high-probability margin-accuracy event,
\[
\widetilde R_0(w)
\le
\widehat R_{\eps_N}(w)
\le
\widetilde R_{2\eps_N}(w)
\qquad
\text{for all }w\in\Delta_K.
\]
On the high-probability VC event,
\[
R_0(w)\le \widetilde R_0(w)+\mathfrak v_Q(\delta),
\qquad
\widetilde R_{2\eps_N}(w)\le R_{2\eps_N}(w)+\mathfrak v_Q(\delta)
\]
for all $w$.  Therefore
\begin{align*}
R(\what_\eps)
&=R_0(\what_\eps)\\
&\le \widetilde R_0(\what_\eps)+\mathfrak v_Q(\delta)\\
&\le \widehat R_{\eps_N}(\what_\eps)+\mathfrak v_Q(\delta)\\
&\le \widehat R_{\eps_N}(\wstar)+\mathfrak v_Q(\delta)\\
&\le \widetilde R_{2\eps_N}(\wstar)+\mathfrak v_Q(\delta)\\
&\le R_{2\eps_N}(\wstar)+2\mathfrak v_Q(\delta).
\end{align*}
Finally,
\[
R_{2\eps_N}(\wstar)-R_0(\wstar)
=
\Pp_q\{0<M_q(\wstar)\le 2\eps_N\}
\le
C_m(2\eps_N)^\beta
\]
by Assumption~\ref{ass:comp-margin}.  Subtract $R(\wstar)=R_0(\wstar)$.
\end{proof}

\begin{remark}[Why the buffer matters]
The theorem avoids any margin assumption at the random, data-dependent $\what_\eps$.  The buffer guarantees that if a training problem appears safely solved empirically, then it is also solved by the true margin, up to the $2\eps_N$ band.  Hence the only finite-$N$ mass we pay for is the $2\eps_N$ band around the fixed comparator $\wstar$.
\end{remark}

\section{Optional fast $Q$-rates from geometric identifiability}

The margin condition alone controls the finite-generation floor.  It does not, by itself, imply a fast $Q$-rate.  A fast $Q$-rate requires an additional condition saying that moving away from $\wstar$ actually increases risk.  This is an identifiability or growth condition.

Define the disagreement probability
\[
D(w,\wstar)
:=
\Pp_q\{\ell_q(w)\ne \ell_q(\wstar)\}.
\]
A common route to fast rates is a Bernstein condition
\[
D(w,\wstar)
\le
B\{R(w)-R(\wstar)\}^{\alpha},
\qquad 0<\alpha\le 1.
\]
Rather than assume this directly, we state primitive geometric conditions that imply it.

\begin{assumption}[Absolute comparator margin]
\label{ass:absolute-margin}
There exist $C_m>0$, $\beta>0$, and $t_0>0$ such that
\[
\Pp_q\{|M_q(\wstar)|\le t\}
\le
C_m t^\beta,
\qquad 0<t\le t_0.
\]
\end{assumption}

\begin{assumption}[Local risk growth and global separation]
\label{ass:growth}
There exist constants $c_g>0$, $\gamma>0$, $r_0>0$, and $c_{\rm sep}>0$ such that
\[
R(w)-R(\wstar)
\ge
c_g\norm{w-\wstar}_1^\gamma
\qquad
\text{whenever }\norm{w-\wstar}_1\le r_0,
\]
and
\[
R(w)-R(\wstar)
\ge c_{\rm sep}
\qquad
\text{whenever }\norm{w-\wstar}_1>r_0.
\]
\end{assumption}

\begin{proposition}[Geometry implies Bernstein]
\label{prop:geom-bernstein}
Under Assumptions~\ref{ass:absolute-margin} and~\ref{ass:growth}, there is a constant $B<\infty$ such that
\[
D(w,\wstar)
\le
B\{R(w)-R(\wstar)\}^{\alpha}
\qquad
\text{for all }w\in\Delta_K,
\]
where
\[
\alpha:=\min\left\{1,\frac{\beta}{\gamma}\right\}.
\]
\end{proposition}

\begin{proof}
First note that $M_q$ is one-Lipschitz in $\ell_1$:
\[
|M_q(w)-M_q(w')|
\le
\norm{w-w'}_1.
\]
Indeed, each $\Delta_{q,a}$ has $\ell_\infty$ norm at most one, and taking the minimum over $a$ preserves the Lipschitz constant.  If $\ell_q(w)\ne\ell_q(\wstar)$, then the signs of $M_q(w)$ and $M_q(\wstar)$ differ, so
\[
|M_q(\wstar)|
\le
|M_q(w)-M_q(\wstar)|
\le
\norm{w-\wstar}_1.
\]
Thus, by Assumption~\ref{ass:absolute-margin},
\[
D(w,\wstar)
\le
C_m\norm{w-\wstar}_1^\beta
\]
whenever $\norm{w-\wstar}_1\le t_0$.  Locally, Assumption~\ref{ass:growth} gives
\[
\norm{w-\wstar}_1^\beta
\le
c_g^{-\beta/\gamma}
\{R(w)-R(\wstar)\}^{\beta/\gamma}.
\]
If $\beta/\gamma>1$, weaken the exponent to one, since $R(w)-R(\wstar)\le 1$ after adjusting constants.  Away from the local ball, $D(w,\wstar)\le1\le c_{\rm sep}^{-\alpha}\{R(w)-R(\wstar)\}^{\alpha}$.  Enlarging constants proves the claim.
\end{proof}

We will use the following standard localized-VC lemma.  It is included as a black-box empirical-process step; the substantive modeling content for this paper is Proposition~\ref{prop:geom-bernstein}, which gives primitive geometric conditions under which the Bernstein premise of the lemma holds.

\begin{lemma}[Localized VC bound for approximate ERM]
\label{lem:localized-vc}
Let $\mathcal L=\{\ell_w:w\in\Delta_K\}$ be a class of $\{0,1\}$-valued losses with minimizer $\wstar$.  Suppose its growth function on $Q$ points is bounded by
\[
\Pi_{\mathcal L}(Q)\le \exp(V_Q),
\qquad
V_Q\asymp (K-1)\log(eAQ),
\]
and suppose the Bernstein condition holds:
\[
\Pp_q\{\ell_q(w)\ne\ell_q(\wstar)\}
\le
B\{R(w)-R(\wstar)\}^{\alpha},
\qquad 0<\alpha\le 1.
\]
If a data-dependent $\widetilde w$ satisfies
\[
\frac1Q\sum_{j=1}^Q\ell_{q_j}(\widetilde w)
\le
\frac1Q\sum_{j=1}^Q\ell_{q_j}(\wstar)+a,
\]
then with probability at least $1-\delta$,
\[
R(\widetilde w)-R(\wstar)
\le
C
\left[
 a
 +
\left(\frac{V_Q+\log(1/\delta)}{Q}\right)^{1/(2-\alpha)}
+
\frac{V_Q+\log(1/\delta)}{Q}
\right],
\]
where $C$ depends only on $B$ and $\alpha$.
\end{lemma}

\begin{proof}[Proof idea]
This is the usual peeling proof for bounded excess losses under a Bernstein condition.  For the excess class $g_w=\ell_w-\ell_{\wstar}$, the variance satisfies $\E g_w^2\le B(\E g_w)^\alpha$.  On each shell $r<\E g_w\le2r$, Bernstein's inequality and the VC growth bound imply a uniform deviation of order
\[
\sqrt{\frac{r^\alpha(V_Q+\log(1/\delta))}{Q}}
+
\frac{V_Q+\log(1/\delta)}{Q}.
\]
Solving the fixed point
\[
r\asymp
\sqrt{\frac{r^\alpha(V_Q+\log(1/\delta))}{Q}}
\]
gives $r\asymp((V_Q+\log(1/\delta))/Q)^{1/(2-\alpha)}$.  The approximate-ERM slack $a$ enters additively.  This is the standard Massart/Tsybakov localized-rate argument, specialized to the finite-growth class above.
\end{proof}

We now state the fast-rate version.

\begin{theorem}[Buffered ERM with optional fast $Q$-rate]
\label{thm:fast}
Assume Assumptions~\ref{ass:absolute-margin} and~\ref{ass:growth}.  Let $\alpha=\min\{1,\beta/\gamma\}$ and let $\what_\eps$ be the buffered ERM from Theorem~\ref{thm:buffered-main}.  Then, with probability at least $1-\delta$,
\[
R(\what_\eps)-R(\wstar)
\le
C
\left(
\frac{(K-1)\log(eAQ)+\log(1/\delta)}{Q}
\right)^{\frac{1}{2-\alpha}}
+
C'
\left(\frac{\log(KAQ/\delta)}{N}\right)^{\beta/2},
\]
where $C,C'$ depend on the margin and growth constants but not on $Q,N,K,A$.
\end{theorem}

\begin{proof}[Proof sketch]
On the margin-accuracy event,
\[
\frac1Q\sum_{j=1}^Q\ell_{q_j}(\what_\eps)
\le
\frac1Q\sum_{j=1}^Q\one\{\widehat M_{q_j}(\what_\eps)\le\eps_N\}
\le
\frac1Q\sum_{j=1}^Q\one\{\widehat M_{q_j}(\wstar)\le\eps_N\}
\le
\frac1Q\sum_{j=1}^Q\one\{M_{q_j}(\wstar)\le2\eps_N\}.
\]
Therefore $\what_\eps$ is an approximate ERM for the clean loss, with empirical excess bounded by
\[
\frac1Q\sum_{j=1}^Q\one\{0<M_{q_j}(\wstar)\le2\eps_N\}.
\]
By the margin condition, this empirical band mass has expectation at most $C_m(2\eps_N)^\beta$ and, by Bernstein's inequality for a fixed Bernoulli variable, is at most a constant multiple of $(2\eps_N)^\beta+\log(1/\delta)/Q$ with high probability.  Proposition~\ref{prop:geom-bernstein} gives the Bernstein condition with exponent $\alpha$.  Applying Lemma~\ref{lem:localized-vc} yields the displayed bound, absorbing the $\log(1/\delta)/Q$ term into the localized complexity term.
\end{proof}

\begin{remark}[Where the exponent comes from]
The exponent is not inherently $\beta/(\beta+1)$.  In this geometry it is
\[
\alpha=\min\{1,\beta/\gamma\},
\]
where $\beta$ measures the thickness of the decision boundary and $\gamma$ measures how risk grows as one moves away from $\wstar$.  The often-appearing value $\alpha=\beta/(\beta+1)$ corresponds to the special local-growth regime $\gamma=\beta+1$.
\end{remark}

\section{Why margin alone cannot give a fast $Q$-rate}

The following example shows why the fast $Q$-rate cannot follow from the geometric margin condition alone, even with a hard margin and no finite-generation noise.

\begin{proposition}[Hard margin but only root-$Q$ distinguishability]
\label{prop:counterexample}
There is a family of two-model, two-answer problems with $N=\infty$ and hard margin at the optimum such that no procedure can uniformly guarantee excess risk $o(Q^{-1/2})$ from $Q$ labeled questions.
\end{proposition}

\begin{proof}
Take $K=2$ and two answers, the gold answer $a^\star$ and a wrong answer $b$.  Write $w=(t,1-t)$.  There are two problem types.

For type $+$, model 1 is always correct and model 2 is always wrong:
\[
p_1(a^\star\mid q)=1,
\qquad
p_2(a^\star\mid q)=0.
\]
Then
\[
\Delta_q=(1,-1),
\qquad
M_q(w)=2t-1.
\]
For type $-$, model 1 is always wrong and model 2 is always correct:
\[
p_1(a^\star\mid q)=0,
\qquad
p_2(a^\star\mid q)=1,
\]
so
\[
\Delta_q=(-1,1),
\qquad
M_q(w)=1-2t.
\]
Thus weights with $t>1/2$ solve exactly the type $+$ problems, while weights with $t<1/2$ solve exactly the type $-$ problems.

Now consider two distributions over problem types:
\[
\D_+:
\Pp(q=+)=\frac12+\frac\eps2,
\qquad
\D_-:
\Pp(q=+)=\frac12-\frac\eps2.
\]
Under $\D_+$, the optimal side is $t>1/2$, with risk $1/2-\eps/2$.  Choosing the wrong side has risk $1/2+\eps/2$, so the excess risk is $\eps$.  Under $\D_-$, the optimal side is reversed.

At an optimum, say $w^\star=(1,0)$ under $\D_+$, the margins satisfy
\[
|M_q(w^\star)|=1
\]
for both problem types.  Hence the margin is hard: there is no near-boundary mass at all.

However, distinguishing $\D_+$ from $\D_-$ using $Q$ training questions is just distinguishing two coins with biases $1/2\pm\eps/2$.  If $\eps=c/\sqrt Q$ for sufficiently small constant $c$, the total variation distance between the $Q$-sample laws remains bounded away from one.  Therefore any procedure chooses the wrong side with constant probability under at least one of $\D_+$ or $\D_-$.  On that event its excess risk is $\eps\asymp Q^{-1/2}$.

Thus even a hard margin does not imply a faster-than-root-$Q$ learning rate.  The missing ingredient is global identifiability: the two separated regions of the simplex have nearly equal risks.
\end{proof}

\section{How to interpret the growth exponent $\gamma$}

The growth condition is not a technical decoration; it describes how the geometry of the benchmark changes when we move the ensemble weights.  Let $r=\norm{w-\wstar}_1$.  Because $M_q$ is Lipschitz in $w$, only problems with
\[
|M_q(\wstar)|\lesssim r
\]
can change their solved/unsolved status.  Under the margin condition, the mass of such problems is at most order $r^\beta$.  This is the disagreement scale:
\[
D(w,\wstar)\lesssim r^\beta.
\]
But the excess risk is not the disagreement mass; it is the net imbalance between newly lost and newly gained questions.

Along a one-dimensional path $w_r$ away from $\wstar$, write
\[
B_+(r):=\Pp\{\text{question is solved by }\wstar\text{ but failed by }w_r\},
\]
\[
B_-(r):=\Pp\{\text{question is failed by }\wstar\text{ but solved by }w_r\}.
\]
Then
\[
D(w_r,\wstar)=B_+(r)+B_-(r),
\qquad
R(w_r)-R(\wstar)=B_+(r)-B_-(r).
\]
The margin exponent controls the sum; the growth exponent controls the difference.

\paragraph{Case A: all crossings hurt.}
If moving away from $\wstar$ mostly turns correct questions into errors and does not rescue comparable failed questions, then
\[
B_+(r)\asymp r^\beta,
\qquad
B_-(r)=o(r^\beta).
\]
Hence
\[
R(w_r)-R(\wstar)\asymp r^\beta,
\]
so $\gamma=\beta$, $\alpha=1$, and the fast $Q$-term is essentially $Q^{-1}$ up to logarithms.  A simple model is a one-sided boundary: all questions near the boundary are correctly solved at $\wstar$ and become wrong when the weight is perturbed.

\paragraph{Case B: first-order cancellation.}
Suppose the total crossing mass is still order $r^\beta$, but losses and gains nearly cancel:
\[
B_+(r)=c r^\beta+c' r^{\beta+1}+o(r^{\beta+1}),
\]
\[
B_-(r)=c r^\beta-c' r^{\beta+1}+o(r^{\beta+1}).
\]
Then
\[
D(w_r,\wstar)\asymp r^\beta,
\qquad
R(w_r)-R(\wstar)\asymp r^{\beta+1}.
\]
Thus $\gamma=\beta+1$, $\alpha=\beta/(\beta+1)$, and the fast $Q$-term becomes
\[
Q^{-(\beta+1)/(\beta+2)}
\]
up to logarithms.  This is the regime in which the previously suggested exponent naturally appears.  It is not a consequence of the margin condition alone; it is a consequence of margin plus first-order gain/loss cancellation.

\paragraph{Case C: higher-order cancellation.}
If the gain/loss imbalance vanishes to higher order,
\[
B_+(r)-B_-(r)\asymp r^{\beta+s},
\qquad s>1,
\]
then $\gamma=\beta+s$ and
\[
\alpha=\frac{\beta}{\beta+s}.
\]
The fast rate becomes slower.  This is a flatter risk landscape: many questions change status near the boundary, but the net risk difference is small.

\paragraph{Case D: separated near-tie between two regions.}
The counterexample in Proposition~\ref{prop:counterexample} is the extreme nonlocal case.  Two separated regions of the simplex have hard margins internally, but their risks differ by only $\eps$.  The benchmark gives no local boundary signal that decides which region is better; one must estimate a global coin bias from $Q$ problems.  This yields the root-$Q$ obstruction.

\section{Recommended theorem hierarchy for the paper}

The cleanest presentation is:

\begin{enumerate}[label=(\roman*)]
\item State the double-randomness setup and the definitions of $p_i$, $\widehat p_i$, $\Delta_{q,a}$, $M_q$, and $\widehat M_q$.

\item Make the buffered ERM theorem, Theorem~\ref{thm:buffered-main}, the main result.  It requires only a comparator margin condition and gives the robust two-term rate
\[
\sqrt{\frac{K\log(AQ)}{Q}}
+
\left(\frac{\log(KAQ)}{N}\right)^{\beta/2}.
\]

\item Include Proposition~\ref{prop:counterexample} immediately after the main theorem to explain why the margin condition alone cannot produce a fast $Q$-rate.

\item Present the optional fast-rate theorem only after introducing geometric identifiability: absolute margin plus risk growth.  Emphasize that the exponent is
\[
\alpha=\min\{1,\beta/\gamma\},
\]
not automatically $\beta/(\beta+1)$.
\end{enumerate}

This makes the paper honest: the genuinely new LLM-specific contribution is the $Q,N$ decomposition and the finite-generation margin floor.  Fast $Q$-rates are available, but only when the risk landscape has additional identifiable geometry.

\part{Adaptive Test-Time Allocation at a Fixed Weight Vector}
\section{Introduction}

\subsection{Setting and motivation}

A test-time ensemble aggregates repeated generations from $K$ language
models: on problem $q$, model $i$ induces an answer distribution
$p_i(\cdot\mid q)$, and a weight vector $\wfix\in\DeltaK$ scores each
candidate answer $a$ by $\langle \wfix,p_\cdot(a\mid q)\rangle$.  The
companion note on buffered ERM studies how to \emph{learn} $\wfix$ from $Q$
labeled questions when each answer distribution is estimated from $N$
generations, and establishes the two-term rate
\[
R(\what)-R(\wstar)
\;\lesssim\;
\sqrt{\frac{K\log(AQ)}{Q}}
\;+\;
\Big(\frac{\log(KAQ)}{N}\Big)^{\beta/2}.
\]
The two terms have very different practical status.  The $Q$-term is
information-theoretic and binds hard on real benchmarks ($Q=30$ for AIME):
no method can learn rich weights from thirty questions.  But the $N$-term is
\emph{compute}: it is the only term the practitioner controls at inference
time, and nothing forces the same $N$ on every problem.  This paper makes
the generation budget the object of study.

Fix the weight vector $\wfix$ (learned, uniform, or best-single --- and note
$K=1$ covers plurality voting / self-consistency for a single model).
The population margin
\[
M_q(\wfix)
=\min_{a\ne a^\star(q)}\big\langle \wfix,\;p_\cdot(a^\star(q)\mid q)-p_\cdot(a\mid q)\big\rangle
\]
measures how safely the infinite-compute ensemble decides problem $q$: it is
the signal-to-noise ratio of the problem.  Estimating the plurality answer
of a problem with margin $m$ from finite generations is a testing problem
whose difficulty scales as $m^{-2}$; a margin condition with exponent
$\beta$ (Assumption~\ref{ass:margin}) says small-margin problems are rare at
a polynomial rate.  The questions we answer:

\begin{enumerate}
\item If every problem receives the same budget $b$, what excess risk over
$R(\wfix)$ is incurred?  (\emph{Uniform allocation} --- Theorem~\ref{thm:uniform};
this is the finite-generation floor restated in budget language.)
\item If budgets may depend on the problem, how much better can one do?
(\emph{Oracle allocation} --- Theorem~\ref{thm:oracle}.)
\item Can the oracle be matched without knowing the margins?
(\emph{Sequential allocation} --- Theorem~\ref{thm:sequential}.)
\item Are these rates optimal?  (\emph{Minimax lower bounds} ---
Theorems~\ref{thm:lower-uniform} and~\ref{thm:lower-adaptive}.)
\end{enumerate}

\subsection{Results}

Write $b$ for the expected number of \emph{rounds} per problem (one round =
one generation from each model; the factor $K$ affects only constants).
Under the margin condition with exponent $\beta\in(0,2)$:
\[
\underbrace{\mathcal E_{\mathrm{unif}}(b)\;\asymp\;b^{-\beta/2}}_{\text{uniform}}
\qquad\text{versus}\qquad
\underbrace{c\,b^{-\frac{\beta}{2-\beta}}
\;\le\;\mathcal E_{\mathrm{adapt}}(b)\;\le\;
C\Big(\frac{\log b}{b}\Big)^{\frac{\beta}{2-\beta}}}_{\text{adaptive}} .
\]
The adaptive exponent $\beta/(2-\beta)$ strictly exceeds the uniform
$\beta/2$ for every $\beta\in(0,2)$; at $\beta=1$ the comparison is $b^{-1}$
versus $b^{-1/2}$.  The upper bound is achieved by a practical rule: sample
in doubling epochs and stop when the empirical margin exceeds an
anytime-valid confidence radius (``stop when the leader's lead is
significant''), with a budget-determined cap.  The lower bounds are minimax
over the class of populations satisfying the margin condition, so the
uniform-versus-adaptive gap is a property of the problem, not of our proofs.
At $\beta=2$ the problem undergoes a phase transition, which
Section~\ref{sec:phase} characterizes completely: for $\beta>2$ the
resolution cost $\mu_2=\E[M^{-2}]$ is finite and adaptive allocation
achieves excess $e^{-\Theta(b)}$ (with a matching exponential lower bound),
while uniform allocation remains polynomial at $\Theta(b^{-\beta/2})$ ---
a polynomial-versus-exponential separation; at $\beta=2$ itself both bounds
are stretched exponentials $e^{-\Theta(\sqrt b)}$, matched in form via a
scale-invariant hard population.  Section~\ref{sec:phase} also resolves how
the theory reads on finite benchmarks, which are hard-margin populations
belonging to the class for every $\beta$ simultaneously.

The mechanism is elementary and worth stating once: at budget $b$, uniform
allocation is bottlenecked by problems at margin $m\asymp b^{-1/2}$, whose
mass is $\asymp b^{-\beta/2}$; it wastes almost all of its budget driving
already-negligible errors on large-margin problems further down.  An
adaptive allocator spends $\asymp m^{-2}$ on each problem at margin $m$ and
can therefore afford to push its resolution cutoff down to
$m_0\asymp b^{-1/(2-\beta)}\ll b^{-1/2}$, paying total mass
$m_0^\beta=b^{-\beta/(2-\beta)}$.

Two further findings.  First, a structural asymmetry between knowing and
learning (Section~\ref{sec:twosided}): the oracle needs only the
\emph{one-sided} margin condition $\Pp\{0<M_q\le t\}\lesssim t^\beta$ ---
it skips hopeless problems for free --- while any sequential learner must
spend budget to discover that a problem with margin just below zero is
hopeless, and correspondingly requires the \emph{two-sided} condition
$\Pp\{|M_q|\le t\}\lesssim t^\beta$.  With an unconstrained negative band the
adaptive rate provably degrades back to the uniform rate.  Second, the
$K=1$ case is a rate theory for \emph{adaptive self-consistency}
(Remark~\ref{rem:selfconsistency}): sample a single model until its leading
answer's lead over the runner-up is statistically significant.  Adaptive
stopping rules of this type are used heuristically (e.g.\ the Bayesian
stopping of the Best-of-$\infty$ paper); our results identify the rates they
achieve and the exponent $\beta$ that governs them.

\subsection{Relation to prior work}

The allocation problem connects three literatures.  (1) \emph{Test-time
scaling}: compute-optimal scaling studies (Snell et al.) demonstrate
empirically that adaptive compute beats uniform compute; we provide a
minimax theory with explicit exponents for the plurality-vote/ensemble
setting.  (2) \emph{Thresholding bandits and best-arm identification}
(Locatelli, Gutzeit and Carpentier; Jamieson et al.'s lil'UCB; Karnin et
al.): our sequential rule is a per-problem sign test of the margin, and our
lower-bound technique (divergence decomposition with adaptive stopping) is
from this toolbox; the objective differs --- population excess risk against
an infinite-compute reference, with the margin distribution playing the
role of the arm-gap distribution --- so the rates and their $\beta$-dependence
are new.  (3) \emph{Active learning under Tsybakov noise} (Castro and
Nowak; Hanneke): the phenomenon ``adaptivity improves the exponent under a
margin condition'' is the same in spirit; the sampling channel here
(multinomial generations against a plurality threshold) and the budget
accounting are specific to test-time ensembling.

\section{Setup}
\label{sec:setup}

Notation follows the companion note.  Problems $q\sim\D$; gold answer
$a^\star(q)$; candidate set $\A(q)$, $|\A(q)|\le A$; models $i=1,\dots,K$
with answer distributions $p_i(\cdot\mid q)$; weights $\wfix\in\DeltaK$
fixed throughout; margins $M_q=M_q(\wfix)\in[-1,1]$ as above; clean risk
$R(\wfix)=\Pp_q\{M_q\le0\}$ (ties are errors).

\paragraph{Rounds and procedures.}
One \emph{round} on problem $q$ draws one generation from each model,
$(Y_{1,r},\dots,Y_{K,r})$, $Y_{i,r}\sim p_i(\cdot\mid q)$, independent
across models and rounds.  After $n$ rounds, $\widehat p^{(n)}_i(a\mid q)$
denotes the empirical answer frequencies,
$\Mhat^{(n)}_q=\min_{a\ne a^\star(q)}\langle\wfix,\widehat
p^{(n)}_\cdot(a^\star(q)\mid q)-\widehat p^{(n)}_\cdot(a\mid q)\rangle$
the empirical margin, and the \emph{plurality output}
$\widehat a^{(n)}=\arg\max_a\langle \wfix,\widehat p^{(n)}_\cdot(a\mid
q)\rangle$ (ties are errors).  A \emph{procedure} decides, after each round
and measurably in the observed generations (plus outside randomization),
whether to request another round, and on stopping outputs an answer; it
observes answer strings only --- neither $a^\star(q)$ nor $M_q$.  A problem
allocated zero rounds outputs an abstention symbol, always an error.  The
spend on $q$ is $n_q$; the budget constraint is
$\E[n_q]\le b$, jointly over $q\sim\D$, the generations, and the
procedure's randomization.  (Strict total budgets over a finite problem set
are recovered by Bernstein's inequality; Remark~\ref{rem:total-budget}.)
The \emph{procedure risk} and \emph{excess} are
\[
R_\pi=\Pp\{\text{output}\ne a^\star(q)\},
\qquad
\mathcal E_\pi=R_\pi-R(\wfix).
\]

\paragraph{One-sidedness of the excess.}
On $\{M_q\le0\}$ the infinite-compute ensemble is already wrong, so finite
sampling can only help there:
\begin{equation}
\label{eq:onesided}
\mathcal E_\pi
\;\le\;
\E\big[\one\{M_q>0\}\;\Pp(\text{output}\ne a^\star\mid q)\big].
\end{equation}
Upper bounds therefore pay only on the positive side.

\begin{assumption}[One-sided margin condition at $\wfix$]
\label{ass:margin}
There exist $C_m>0$, $\beta>0$, $t_0\in(0,1]$ with
$\Pp_q\{0<M_q(\wfix)\le t\}\le C_m t^\beta$ for all $0<t\le t_0$.
\end{assumption}

\begin{assumption}[Two-sided margin condition at $\wfix$]
\label{ass:two-sided}
There exist $C_m>0$, $\beta>0$, $t_0\in(0,1]$ with
$\Pp_q\{|M_q(\wfix)|\le t\}\le C_m t^\beta$ for all $0<t\le t_0$.
\end{assumption}

Assumption~\ref{ass:two-sided} implies Assumption~\ref{ass:margin} and
additionally forbids an atom at $M_q=0$.  The oracle results use only
Assumption~\ref{ass:margin}; the sequential learner uses
Assumption~\ref{ass:two-sided}, and Section~\ref{sec:twosided} shows this is
not an artifact.  On finite benchmarks the comparator should be chosen with
healthy margins --- among the risk minimizers, tie-breaking can manufacture
spurious atoms at zero (empirically: a sign-MILP parks problems at
$M\approx10^{-4}$, its numerical buffer), and a max-margin selection removes
them.

\section{The cost of a margin}

\begin{lemma}[Per-problem sign error]
\label{lem:perproblem}
Fix $q$ with $m:=M_q(\wfix)>0$.  After $n\ge1$ rounds,
$\Pp\{\widehat a^{(n)}\ne a^\star(q)\}\le(A-1)\,e^{-nm^2/2}$.
\end{lemma}

\begin{proof}
For a wrong answer $a$, the per-round score differences
$X^{(a)}_r=\sum_i\wfix_i(\one\{Y_{i,r}=a^\star(q)\}-\one\{Y_{i,r}=a\})\in[-1,1]$
are i.i.d.\ with mean $\langle\wfix,\Delta_{q,a}\rangle\ge m$, where
$\Delta_{q,a}=p_\cdot(a^\star(q)\mid q)-p_\cdot(a\mid q)$, and their average
is the empirical gold-versus-$a$ gap.  The output is wrong only if some
wrong answer's empirical gap is $\le0$, a deviation of at least $m$ below
the mean; Hoeffding (range $2$) bounds each such event by $e^{-nm^2/2}$, and
a union bound over the $\le A-1$ wrong answers finishes.
\end{proof}

Thus driving the error on a margin-$m$ problem below $\delta$ costs
$n(m,\delta)=\lceil(2/m^2)\log((A-1)/\delta)\rceil$ rounds.  This $m^{-2}$
law is the engine of all upper bounds, and Section~\ref{sec:lower} shows it
is information-theoretically exact.  Empirically it is directly visible in
real LLM generations: in our AIME2024 resampling study, problems with
$|M|\ge0.15$ are decided by $N'\approx10$ generations while near-tie
problems keep flipping at $N'=100$, exactly the $e^{-cNM^2}$ stratification.

\section{Uniform allocation}

\begin{theorem}[Uniform allocation]
\label{thm:uniform}
Under Assumption~\ref{ass:margin}, the rule that spends $n_q\equiv b\ge1$
rounds on every problem and outputs the plurality answer satisfies
\[
\mathcal E_{\mathrm{unif}}(b)
\ \le\
(A-1)\Big[
C_m\,2^{\beta/2}\,\Gamma\big(\tfrac\beta2+1\big)\,b^{-\beta/2}
+e^{-t_0^2b/2}
\Big].
\]
\end{theorem}

\begin{proof}
By \eqref{eq:onesided} and Lemma~\ref{lem:perproblem},
$\mathcal E_{\mathrm{unif}}\le(A-1)\,\E[\one\{M_q>0\}e^{-bM_q^2/2}]$.  The
layer-cake identity $e^{-bm^2/2}=\int_m^\infty bt\,e^{-bt^2/2}\,dt$ and
Fubini give
\[
\E\big[\one\{M_q>0\}e^{-bM_q^2/2}\big]
=\int_0^\infty b\,t\,e^{-bt^2/2}\,\Pp\{0<M_q\le t\}\,dt .
\]
On $(t_0,\infty)$ bound the probability by $1$, contributing
$e^{-bt_0^2/2}$.  On $(0,t_0]$ apply the margin condition and substitute
$s=t\sqrt b$:
$\int_0^{t_0}bt\,e^{-bt^2/2}C_mt^\beta dt
\le C_m b^{-\beta/2}\int_0^\infty s^{\beta+1}e^{-s^2/2}ds
=C_m2^{\beta/2}\Gamma(\tfrac\beta2+1)\,b^{-\beta/2}$.
\end{proof}

The substitution $s=t\sqrt b$ displays the mechanism: the polynomial
aggregate rate is the margin-distribution average of per-problem
exponentials, dominated by problems at margin $\asymp b^{-1/2}$.  With
$b=N$ this is precisely the finite-generation floor of the companion note,
now read as the \emph{fixed-allocation baseline} of a design problem.

\section{Oracle allocation}

The oracle sees $M_q$ and allocates, for parameters $m_0\in(0,t_0]$,
$\lambda\ge0$,
\begin{equation}
\label{eq:oracle-alloc}
n_{\mathrm{or}}(q)
=
\Big\lceil \tfrac{2(\lambda+\log(A-1))}{M_q^2}\Big\rceil
\,\one\{M_q\ge m_0\}
\end{equation}
(abandoning sub-cutoff and hopeless problems --- abandoning $M_q\le0$
problems is free, since the clean rule errs there too).

\begin{theorem}[Oracle allocation, $\beta<2$]
\label{thm:oracle}
Under Assumption~\ref{ass:margin} with $\beta\in(0,2)$, the allocation
\eqref{eq:oracle-alloc} satisfies
\[
\mathcal E_{\mathrm{or}}
\le
C_m m_0^\beta+e^{-\lambda},
\qquad
\E[n_{\mathrm{or}}(q)]
\le
2\big(\lambda+\log(A-1)\big)
\Big[\tfrac{2C_m}{2-\beta}m_0^{\beta-2}+C_mt_0^{\beta-2}+t_0^{-2}\Big]+1 .
\]
\end{theorem}

\begin{proof}
\emph{Excess:} the abandoned band $\{0<M_q<m_0\}$ has mass $\le C_mm_0^\beta$;
on $\{M_q\ge m_0\}$, Lemma~\ref{lem:perproblem} with
$n\ge2(\lambda+\log(A-1))/M_q^2$ gives conditional error $\le e^{-\lambda}$.
Combine via \eqref{eq:onesided}.

\emph{Budget:} with $F(t)=\Pp\{0<M_q\le t\}\le C_mt^\beta$ ($t\le t_0$),
$\E[n_{\mathrm{or}}]\le2(\lambda+\log(A-1))\,\E[M_q^{-2}\one\{M_q\ge
m_0\}]+1$.  Split at $t_0$; on $(t_0,1]$ bound $M_q^{-2}\le t_0^{-2}$.  On
$[m_0,t_0]$ integrate by parts:
\[
\int_{[m_0,t_0]}m^{-2}dF(m)
=\big[F(m)m^{-2}\big]_{m_0^-}^{t_0}+2\int_{m_0}^{t_0}F(m)m^{-3}dm
\le C_mt_0^{\beta-2}+\tfrac{2C_m}{2-\beta}\,m_0^{\beta-2},
\]
dropping the nonpositive lower boundary term and using $\beta<2$ (the
divergence of $\int m^{\beta-3}dm$ at the lower limit is exactly why a
cutoff is necessary).
\end{proof}

\begin{corollary}[Oracle rate]
\label{cor:oracle-rate}
With $m_0=(\log b/b)^{1/(2-\beta)}\wedge t_0$ and
$\lambda=\tfrac{\beta}{2-\beta}\log b$: expected budget $O(b)$ and
\[
\mathcal E_{\mathrm{or}}(b)
\ \lesssim\
\Big(\frac{\log b}{b}\Big)^{\beta/(2-\beta)} .
\]
\end{corollary}

\begin{remark}[Why $\beta<2$]
\label{rem:phase}
The constraint $\beta<2$ enters through the divergence of
$\int m^{\beta-3}\,dm$ at the origin: for $\beta<2$ the expected resolution
cost $\E[M^{-2}]$ is infinite and a cutoff is unavoidable.  For $\beta\ge2$
the cost is (log-)integrable and the rates change character entirely ---
adaptive allocation becomes exponentially efficient.
Section~\ref{sec:phase} develops this phase transition, with matching lower
bounds.
\end{remark}

\section{Sequential allocation without an oracle}
\label{sec:sequential}

Parameters: cap $\nmax\ge2$, failure level $\delta\in(0,1)$; set
$J=\lceil\log_2\nmax\rceil$, $L=\log\frac{2(A-1)J}{\delta}$,
$\eps(n)=\sqrt{2L/n}$.

\medskip
\noindent\textbf{Sequential rule} (independently on each problem):
for epochs $j=1,\dots,J$ with sizes $n_j=\min\{2^j,\nmax\}$, sample to $n_j$
total rounds and compute $\Mhat^{(n_j)}_q$; \textbf{stop} and output the
plurality answer if $|\Mhat^{(n_j)}_q|>\eps(n_j)$, or upon reaching the cap.

\medskip
The rule tests the sign of $M_q$: keep sampling while the interval
$\Mhat\pm\eps(n)$ straddles zero.  It knows neither $M_q$, nor $\beta$, nor
$\D$.

\begin{lemma}[Anytime accuracy]
\label{lem:anytime}
For each $q$, the event
$G_q=\{|\Mhat^{(n_j)}_q-M_q|\le\eps(n_j)\ \forall j\le J\}$ has
$\Pp(G_q^c\mid q)\le\delta$.
\end{lemma}

\begin{proof}
For fixed $j$ and wrong answer $a$, two-sided Hoeffding on the first $n_j$
per-round differences gives deviation probability $\le2e^{-n_j\eps(n_j)^2/2}
=2e^{-L}=\delta/((A-1)J)$; union over $a$ and the $J$ checkpoints; the
empirical margin (a minimum over $a$) moves by at most the largest pairwise
deviation.  Sample reuse across epochs is immaterial: each checkpoint
constrains the average of the first $n_j$ rounds, and the union bound is
over fixed checkpoints.
\end{proof}

\begin{lemma}[Behavior on the good event]
\label{lem:behavior}
Set $\bar m=2\eps(\nmax)$ and $\bar m'=4\sqrt{L/\nmax}$.  On $G_q$:
(i) any stop by threshold has the correct sign, and a positive-sign stop
outputs $a^\star(q)$;
(ii) if $|M_q|>0$ then $n_q\le\min\{16L/M_q^2,\nmax\}$, and every problem
with $|M_q|\ge\bar m$ decides before the cap;
(iii) reaching the cap undecided implies $|M_q|\le\bar m$.
\end{lemma}

\begin{proof}
(i) $\Mhat^{(n_j)}_q>\eps(n_j)$ forces $M_q>0$ on $G_q$, and
$\Mhat^{(n_j)}_q>0$ means the gold answer strictly leads empirically; the
negative case is symmetric.  (ii) Let $m=|M_q|$ and $j^\ast$ the first epoch
with $\eps(n_{j^\ast})<m/2$; on $G_q$,
$|\Mhat^{(n_{j^\ast})}_q|\ge m-\eps(n_{j^\ast})>\eps(n_{j^\ast})$, so the
rule stops by epoch $j^\ast$; since $\eps(n)<m/2\iff n>8L/m^2$ and epochs
double, $n_{j^\ast}\le16L/m^2$.  For $m\ge\bar m$, $8L/m^2\le\nmax$.
(iii) is the contrapositive.
\end{proof}

\begin{theorem}[Sequential allocation]
\label{thm:sequential}
Under Assumption~\ref{ass:two-sided} with $\beta\in(0,2)$, if
$\bar m'=4\sqrt{L/\nmax}\le t_0$ then
\[
\mathcal E_{\mathrm{seq}}
\le
\delta+C_m\bar m^\beta,
\qquad
\E[n_q]
\le
16L\Big[\tfrac{2C_m}{2-\beta}(\bar m')^{\beta-2}
+C_m(\bar m')^{\beta-2}+C_mt_0^{\beta-2}+t_0^{-2}\Big]
+\delta\nmax+2 .
\]
\end{theorem}

\begin{proof}
\emph{Excess:} on $G_q^c$ pay $\delta$ (Lemma~\ref{lem:anytime}).  On $G_q$:
problems with $M_q>\bar m$ stop with the positive decision and are correct
(Lemma~\ref{lem:behavior}(i),(ii)); the band $0<M_q\le\bar m$ pays its mass
$\le C_m\bar m^\beta$ (Assumption~\ref{ass:two-sided}, $\bar m\le t_0$);
problems with $M_q\le0$ contribute nonpositively, by \eqref{eq:onesided}.

\emph{Budget:} on $G_q$, $n_q\le\min\{16L/M_q^2,\nmax\}$; on $G_q^c$,
$n_q\le\nmax$, contributing $\delta\nmax$.  Split the minimum at
$|M_q|\lessgtr\bar m'$ (where $16L/(\bar m')^2=\nmax$):
$\nmax\Pp\{|M_q|\le\bar m'\}\le16LC_m(\bar m')^{\beta-2}$ by
Assumption~\ref{ass:two-sided}, and
$\E[M_q^{-2}\one\{|M_q|>\bar m'\}]$ is bounded by the same integration by
parts as in Theorem~\ref{thm:oracle} with the two-sided
$F(t)=\Pp\{0<|M_q|\le t\}$.  The $+2$ absorbs epoch rounding.
\end{proof}

\begin{corollary}[Sequential rate]
\label{cor:seq-rate}
Choosing $\nmax=\lceil b^{2/(2-\beta)}\rceil$, $\delta=1/\nmax$ yields
expected spend $O(b(\log b)^{\beta/2})$ and excess
$O((\log b)^{\beta/2}\,b^{-\beta/(2-\beta)})$; re-normalizing the budget,
with expected spend $\le b$,
\[
\mathcal E_{\mathrm{seq}}(b)
\ \lesssim\
\Big(\frac{\log b}{b}\Big)^{\beta/(2-\beta)},
\]
matching the oracle rate of Corollary~\ref{cor:oracle-rate} including the
power of the logarithm.
\end{corollary}

\begin{remark}[Adaptive self-consistency: the case $K=1$]
\label{rem:selfconsistency}
With one model the rule reads: \emph{sample until the leading answer's lead
over the runner-up exceeds the confidence radius, or a cap is hit; output
the leader}.  Corollary~\ref{cor:seq-rate} is thus a rate theory for
adaptive self-consistency, with $\beta$ measuring the prevalence of
questions on which the model is nearly torn between answers.  Nothing in
the proofs uses $K>1$.
\end{remark}

\begin{remark}[Strict total budgets]
\label{rem:total-budget}
Over $T$ i.i.d.\ problems the total spend concentrates: each $n_{q_t}$ is
bounded by $\nmax$, so Bernstein gives
$\sum_tn_{q_t}\le2T\E[n_q]$ with probability
$1-e^{-cT\E[n_q]/\nmax}$, which is high as soon as
$T\gg b^{\beta/(2-\beta)}$ (up to logs).  A hard global stop at $2Tb$ then
never triggers on this event, leaving the risk bound unchanged.
\end{remark}

\section{Minimax lower bounds}
\label{sec:lower}

Both rates are optimal over the class defined by the margin condition.  Let
$\class(\beta,C_m,t_0)$ be the set of populations satisfying
Assumption~\ref{ass:margin} at $\wfix$.  The hard populations below use
two-answer problems with all $K$ models identical,
$p_i(g\mid q)=\frac{1+m}{2}$ for a gold label $g$ drawn uniformly from the
two answers, independently per problem; then $M_q(\wfix)=m$ for every
$\wfix$, one round is informationally one $\Ber(\frac{1+m}{2})$ observation,
all margins are positive (so $R(\wfix)=0$ and excess $=$ error probability),
and the constructions lie in $\class$ --- and in its two-sided refinement ---
for every $K$, $A\ge2$, and every fixed $\wfix$.  The label randomization is
what makes the lower bound binding for arbitrary decision rules: the data
identify the gold answer only through which answer leads.

\begin{lemma}[Two-point bound with adaptive stopping]
\label{lem:lecam}
Fix margin $m\le\frac12$ and $g$ uniform on the two answers.  For any
stopping rule $\tau$ and output $\widehat a$,
\[
\Pp\{\widehat a\ne g\}
\ \ge\
\frac{1-\TV(P_{\mathsf u},P_{\mathsf v})}{2},
\qquad
\TV(P_{\mathsf u},P_{\mathsf v})
\le
m\sqrt{2\min\{\E_{\mathsf u}[\tau],\E_{\mathsf v}[\tau]\}},
\]
where $P_{\mathsf u},P_{\mathsf v}$ are the joint laws of
$(\tau,\text{observations},\widehat a)$ under the two labels.  In
particular, $\E[\tau]\le\frac{1}{8m^2}$ forces error probability
$\ge\frac14$.
\end{lemma}

\begin{proof}
Le Cam's two-point inequality gives the first display.  For the second:
under the two labels the rounds are i.i.d.\ $\Ber(\frac{1\pm m}{2})$, with
$\KL(\Ber(\tfrac{1+m}{2})\Vert\Ber(\tfrac{1-m}{2}))
=m\log\frac{1+m}{1-m}\le4m^2$ on $(0,\frac12]$.  The divergence
decomposition for adaptively stopped i.i.d.\ sequences (Wald's identity on
the log-likelihood ratio) gives
$\KL(P_{\mathsf u}\Vert P_{\mathsf v})=\E_{\mathsf u}[\tau]\,
\KL(\Ber(\tfrac{1+m}{2})\Vert\Ber(\tfrac{1-m}{2}))\le4m^2\E_{\mathsf
u}[\tau]$ (the output adds nothing, by data processing); Pinsker converts
to $\TV$; the symmetric argument gives the other branch of the minimum.
\end{proof}

\begin{theorem}[Uniform allocation is stuck at $b^{-\beta/2}$]
\label{thm:lower-uniform}
For every $\beta>0$ and $b\ge t_0^{-2}/8$ there is
$\D_{\mathrm u}\in\class(\beta,C_m,t_0)$ such that every uniform-allocation
procedure with budget $b$ has
$\mathcal E_\pi(\D_{\mathrm u})\ge\frac{C_m}{4}(8b)^{-\beta/2}\wedge\frac18$.
\end{theorem}

\begin{proof}
Let $\D_{\mathrm u}$ draw $m$ with CDF $\frac{C_m}{2}t^\beta$ on
$(0,t_1]$, $t_1=\min\{t_0,(2/C_m)^{1/\beta}\}$, remaining mass at $m=1$.
Problems with $m\le m_b:=(8b)^{-1/2}$ have, at deterministic spend $b$,
$\TV\le m_b\sqrt{2b}=\frac12$ (Lemma~\ref{lem:lecam}), hence conditional
error $\ge\frac14$; their mass is $\frac{C_m}{2}m_b^\beta$ when $m_b\le
t_1$, else $\ge\frac12$.
\end{proof}

\begin{theorem}[No strategy beats $b^{-\beta/(2-\beta)}$]
\label{thm:lower-adaptive}
Let $\beta\in(0,2)$.  There are $c_1,b_0>0$ (depending on
$\beta,C_m,t_0$) such that for every $b\ge b_0$ there is
$\D_b\in\class(\beta,C_m,t_0)$ with: every procedure with expected
per-problem spend $\le b$ has
$\mathcal E_\pi(\D_b)\ge c_1\,b^{-\beta/(2-\beta)}$.
\end{theorem}

\begin{proof}
Set $m_0=(\frac{C_m}{64b})^{1/(2-\beta)}$ and $\mu=\frac{C_m}{2}m_0^\beta$;
let $\D_b$ put mass $\mu$ on \emph{atom} problems with margin $m_0$ and
mass $1-\mu$ on deterministic problems ($m=1$), labels uniform.  ($b_0$
ensures $m_0\le t_0\wedge\frac12$, $\mu\le\frac12$; membership in $\class$
is immediate, and the population also satisfies the two-sided condition and
the matching lower mass bound, so it lies in every version of the class.)
By Markov, the conditional expected spend on atoms is $T\le b/\mu$, so
$\min\{\E_{\mathsf u}\tau,\E_{\mathsf v}\tau\}\le2b/\mu$, and by
Lemma~\ref{lem:lecam},
\[
\TV\le m_0\sqrt{\tfrac{4b}{\mu}}
=\sqrt{\tfrac{8\,b\,m_0^{2-\beta}}{C_m}}
=\sqrt{\tfrac{8}{64}}\le\tfrac12
\quad\Longrightarrow\quad
\mathcal E_\pi\ge\frac\mu4=\frac{C_m}{8}
\Big(\frac{C_m}{64}\Big)^{\beta/(2-\beta)}b^{-\beta/(2-\beta)} .
\]
Cross-problem adaptivity changes nothing: labels are independent across
problems, so other problems' data are independent of an atom's label, and
only the marginal expected spend entered.
\end{proof}

\begin{remark}[The rates are pinned]
For $\beta\in(0,2)$, over $\class(\beta,C_m,t_0)$:
uniform $\asymp b^{-\beta/2}$; adaptive between $c\,b^{-\beta/(2-\beta)}$
and $C(\log b/b)^{\beta/(2-\beta)}$.  The uniform-versus-adaptive exponent
gap is real.  For $\beta\ge2$ no polynomial lower bound is possible: the
adversary's plantable mass $\mu\asymp m_0^\beta$ times the resolution cost
$m_0^{-2}$ vanishes, so the spend-accounting step yields nothing ---
consistent with the exponential upper bounds beyond the transition.
Anti-concentration is built into the constructions (Bernoulli rounds have
variance $\ge\frac34$ near the boundary); a lower bound at a \emph{fixed}
population would need it as an assumption.
\end{remark}

\section{The full range of $\beta$: a phase transition at $\beta=2$}
\label{sec:phase}

The rates of Sections 4--7 concern $\beta\in(0,2)$.  This section completes
the picture.  The organizing quantity is the expected resolution cost
\[
\mu_2:=\E\big[M_q^{-2}\one\{M_q>0\}\big],
\]
which is finite precisely when the margin condition holds with some
$\beta>2$ (integration by parts; the boundary term $F(m)m^{-2}\le
C_mm^{\beta-2}$ vanishes at $0^+$ exactly then).  The results of this
section, in summary:

\begin{center}
\begin{tabular}{lll}
\hline
regime & uniform & adaptive \\
\hline
$\beta\in(0,2)$ & $\Theta(b^{-\beta/2})$ &
$\Theta(b^{-\beta/(2-\beta)})$ up to $\log$ \\
$\beta=2$ & $\Theta(b^{-1})$ & $e^{-\Theta(\sqrt b)}$ \\
$\beta>2$ & $\Theta(b^{-\beta/2})$ & $e^{-\Theta(b)}$ \\
hard margin ($\mmin>0$) & $\le e^{-b\mmin^2/2}$ &
$\le e^{-(b-1)/(2\mu_2)}$ \\
\hline
\end{tabular}
\end{center}

\noindent
Uniform allocation has \emph{no} transition: Theorem~\ref{thm:lower-uniform}
holds for every $\beta>0$, so on continuum populations the uniform rate
merely steepens as $b^{-\beta/2}$.  The transition is intrinsically about
adaptivity --- only an allocator that spends $\asymp M_q^{-2}$ where needed,
rather than $b$ everywhere, can exploit the integrability of the resolution
cost.  For $\beta>2$ the uniform-versus-adaptive separation is therefore
polynomial versus exponential.

In the exponential regime, Pinsker's inequality (the route to
Lemma~\ref{lem:lecam}) is vacuous; the lower bounds instead use the
Bretagnolle--Huber inequality, which combined with the divergence
decomposition for stopped sequences gives, in the setting of
Lemma~\ref{lem:lecam},
\begin{equation}
\label{eq:bh}
\Pp\{\widehat a\ne g\}
\ \ge\
\tfrac14\,e^{-8m^2T},
\qquad
T=\text{label-averaged expected spend.}
\end{equation}

\begin{theorem}[$\beta>2$: exponential rates, matched]
\label{thm:beta-gt2}
Let Assumption~\ref{ass:margin} hold with $\beta>2$, so
$\mu_2\le C_mt_0^{\beta-2}(1+\tfrac2{\beta-2})+t_0^{-2}<\infty$.
\begin{enumerate}
\item[(i)] The cutoff-free allocation
$n(q)=\lceil2(\lambda+\log(A-1))/M_q^2\rceil\one\{M_q>0\}$ achieves, at
expected spend $b$,
\[
\mathcal E_{\mathrm{or}}(b)\ \le\ (A-1)\,e^{-(b-1)/(2\mu_2)} .
\]
\item[(ii)] Conversely, with $m_*=t_0$ and
$\mu_0=\min\{C_mt_0^\beta,\frac12\}$, the population placing mass $\mu_0$
at margin $m_*$ (gold labels uniform) and the rest at margin $1$ lies in
$\class(\beta,C_m,t_0)$, and every procedure with expected spend $\le b$
suffers
\[
\mathcal E_\pi\ \ge\ \frac{\mu_0}{4}\,e^{-8m_*^2b/\mu_0} .
\]
When the atom dominates $\mu_2$ (i.e.\ $\mu_0m_*^{-2}\ge1$), the two
exponential rate constants agree up to a universal factor.
\end{enumerate}
\end{theorem}

\begin{proof}
(i): as in Theorem~\ref{thm:oracle} with $m_0=0$ now admissible: every
positive-margin problem is driven to conditional error $e^{-\lambda}$ at
expected cost $2(\lambda+\log(A-1))\mu_2+1$; invert the budget.  (ii):
class membership as in Theorem~\ref{thm:lower-adaptive}; Markov gives
label-averaged spend $T\le b/\mu_0$ on the atom; apply \eqref{eq:bh} and
multiply by the atom mass.  For the constants comparison,
$\mu_2=\mu_0m_*^{-2}+(1-\mu_0)\le2\mu_0m_*^{-2}$ when
$\mu_0m_*^{-2}\ge1$, so $8m_*^2/\mu_0\le16/\mu_2$.
\end{proof}

\begin{theorem}[$\beta=2$: stretched exponentials, matched in form]
\label{thm:beta2}
Let $\beta=2$.
\begin{enumerate}
\item[(i)] Under Assumption~\ref{ass:margin}, oracle allocation with
cutoff $m_0$ and $\lambda=2\log(t_0/m_0)$ spends
$O\big((\log(t_0/m_0))^2\big)$ in expectation (the budget integral is now
logarithmic), so at expected spend $b$,
$\mathcal E_{\mathrm{or}}(b)\le Ce^{-c\sqrt b}$.
\item[(ii)] Conversely, let $c_0=\frac34C_mt_0^2\wedge\frac13$ and let
$\D_b$ be the \emph{geometric ladder}: margins $m_j=t_02^{-j}$ with masses
$\mu_j=c_04^{-j}$, $j\le J$, gold labels uniform, remaining mass at margin
$1$.  Then $\D_b\in\class(2,C_m,t_0)$, and with
$J\asymp\sqrt b$, every procedure with expected spend $\le b$ suffers
$\mathcal E_\pi\ge\frac{c_0}{4}e^{-C_\star\sqrt b}$, with $C_\star$
depending only on $(C_m,t_0)$.
\end{enumerate}
\end{theorem}

\begin{proof}
(i): at $\beta=2$,
$\int_{m_0}^{t_0}F(m)m^{-3}dm\le C_m\log(t_0/m_0)$ in the budget bound of
Theorem~\ref{thm:oracle}; with $\lambda=2\log(t_0/m_0)$ the spend is
$\lesssim(\log(t_0/m_0))^2$ and the excess is $\lesssim(m_0/t_0)^2$; choose
$\log(t_0/m_0)\asymp\sqrt b$.  (ii): membership follows since
$\Pp\{0<M\le t\}\le\frac43c_04^{-k}\le C_mt^2$ for $t\in[m_k,m_{k-1})$.
The ladder is \emph{scale-invariant}: $m_j^2/\mu_j=t_0^2/c_0$ for every
$j$ --- every scale is equally expensive to clear, and there are $J$ of
them.  Let $T_j$ be the label-averaged expected spend at scale $j$; the
budget forces some $j^\ast$ with $\mu_{j^\ast}T_{j^\ast}\le b/J$, whence by
\eqref{eq:bh} and scale invariance the conditional error at scale $j^\ast$
is $\ge\frac14e^{-(8t_0^2/c_0)\,b/J}$, and
$\mathcal E_\pi\ge\mu_{j^\ast}\cdot\frac14e^{-(8t_0^2/c_0)b/J}
\ge\frac{c_0}4e^{-(8t_0^2/c_0)b/J-2J\log2}$.  Optimizing $J\asymp\sqrt b$
balances the exponents at $C_\star\sqrt b$.
\end{proof}

The scale-invariant ladder is the structural signature of criticality: it
is the discrete version of the margin density that makes $\mu_2$ diverge
logarithmically.  We find it noteworthy that the pigeonhole argument
closes the critical case with no logarithmic gap in the exponent's power
of $b$.

\subsection{Hard margins: how the theory reads on a finite benchmark}
\label{sec:hard-margin}

Every real benchmark is a finite atomic population.  Such a population has
a minimum positive margin $\mmin>0$, hence satisfies the margin condition
\emph{vacuously for every $\beta$ simultaneously} (take $t_0=\mmin/2$; the
near-boundary band is empty).  In particular, no exponent $\beta$ is
identified by class membership on a finite benchmark --- consistent with our
companion validation study, which fits only \emph{band-limited effective}
exponents.  The correct asymptotic statement is exponential:

\begin{proposition}[Hard-margin populations]
\label{prop:hard-margin}
Let $\D$ be finitely supported with
$\mmin=\min\{M_q:M_q>0\}>0$, $p_+=\Pp\{M_q>0\}$.  Then
\[
\mathcal E_{\mathrm{unif}}(b)\le(A-1)\,p_+\,e^{-b\mmin^2/2},
\qquad
\mathcal E_{\mathrm{adapt}}(b)\le(A-1)\,e^{-(b-1)/(2\mu_2)},
\]
and the ratio of the exponential rate constants is
\[
\frac{1/(2\mu_2)}{\mmin^2/2}
=\frac{1}{\E\big[(\mmin/M_q)^2\one\{M_q>0\}\big]}\ \ge\ 1 :
\]
an \emph{effective number of hardest problems}.  Uniform allocation's rate
constant is set by the single hardest problem; adaptive allocation's by the
average resolution cost.  If one problem among $Q_0$ (uniform empirical
measure) is uniquely hardest, the ratio is $\approx Q_0$.
\end{proposition}

\begin{remark}[The polynomial regime is the transient --- and the operating
regime]
\label{rem:transient}
The exponential asymptotics switch on at the crossover budget
$b\approx\mmin^{-2}$, where the hardest problem resolves.  Below it, the
uniform excess $\sum_qP(q)e^{-bM_q^2/2}$ is a sum of exponentials whose
rate constants fill a band, and over the window in which those scales are
extinguished one by one it tracks a power law --- whose log--log slope is
precisely the band-limited effective $\beta$ measured by our resampling
experiments (there, budgets $N'\le100$ sit below the crossover).  Thus:
polynomial rates with the measured effective $\beta$ describe the operating
regime $b\lesssim\mmin^{-2}$ of a finite benchmark; beyond the crossover,
decay is exponential and adaptivity buys the rate-constant factor of
Proposition~\ref{prop:hard-margin}.  For a continuum margin distribution
with $\Pp\{0<M\le t\}\asymp t^\beta$ near zero, the ``transient'' is exact
and permanent --- that is the theory of Sections 4--7.
\end{remark}

\begin{remark}[Relation to bandit hardness]
$\mu_2$ is the population analog of the hardness quantity
$H=\sum_i\Delta_i^{-2}$ of best-arm identification and thresholding
bandits, where exponential error decay likewise has rate constant $1/H$.
The margin condition is then a tail condition on the arm-gap distribution
determining whether $H$ is finite ($\beta>2$) or infinite ($\beta<2$), with
the allocation rates changing character exactly at the integrability
threshold.  We are not aware of a prior statement of this transition in an
allocation-across-problems setting; the closest analog is the exponential
speedup of active learning under bounded (Massart) noise relative to
polynomial rates under Tsybakov noise.
\end{remark}

\section{Knowing versus learning: the two-sided condition}
\label{sec:twosided}

The oracle results use the one-sided margin condition; the sequential
learner uses the two-sided one.  This asymmetry is intrinsic, not an
artifact of the proofs.  A problem with margin slightly below zero
contributes nothing to the excess --- it is lost regardless --- but a learner
cannot know this without paying for it: deciding the sign at $|M|=m$ costs
$\asymp m^{-2}$ rounds (Lemma~\ref{lem:lecam}), and the cost is symmetric in
the sign.  Concretely, augment the population of
Theorem~\ref{thm:lower-adaptive} with mass $\nu$ at margin $-m_0$ ---
permitted by the one-sided class, which does not constrain the negative
side.  A procedure achieving the adaptive excess rate must distinguish
$\pm m_0$ on the positive atoms, and the sampling laws at $\pm m_0$ are
mutually indistinguishable before $\asymp m_0^{-2}$ rounds; symmetry then
forces expected spend $\gtrsim\nu m_0^{-2}$ on the negative atoms.  For
constant $\nu$ this forces $b\gtrsim m_0^{-2}$, collapsing the achievable
cutoff to $m_0\asymp b^{-1/2}$ and the excess to the \emph{uniform} rate
$b^{-\beta/2}$.  (We state this as a sketch; the full two-hypothesis
argument is routine but notationally heavy.)  In words:

\begin{center}
\emph{One-sided margin control is what an oracle needs;\\
two-sided margin control is what a learner's budget buys.}
\end{center}

This retroactively motivates the absolute-margin condition appearing in the
fast-rate analysis of the companion note, and it has a practical echo: weight
vectors that park problems near the zero boundary (as raw sign-MILP
solutions do) are budget poison for adaptive sampling --- yet another reason,
beyond the estimation pathologies documented in the companion note, to
prefer buffered/max-margin weight selection.

\section{Composing with weight learning}
\label{sec:compose}

The natural end-to-end protocol is two-stage: learn $\what$ on $Q$ labeled
problems with $N$ generations each (buffered ERM, companion note), then
deploy $\what$ with sequential allocation at budget $b$ on test problems.
If the margin condition holds at $\wstar$ (two-sided, with the max-margin
comparator convention) and $\norm{\what-\wstar}_1$ is small, margins move by
at most $\norm{\what-\wstar}_1$ (the margin is $1$-Lipschitz in $\ell_1$),
so the margin condition transfers to $\what$ with an adjusted band; the
resulting excess is bounded by the sum of the learning term and the
allocation term,
\[
R_{\mathrm{total}}
\;\lesssim\;
R(\wstar)
+\underbrace{\sqrt{\tfrac{K\log(AQ)}{Q}}
+\Big(\tfrac{\log(KAQ)}{N}\Big)^{\beta/2}}_{\text{learning }\what}
+\underbrace{\Big(\tfrac{\log b}{b}\Big)^{\beta/(2-\beta)}}_{\text{allocating at }\what}.
\]
We state this informally; the bookkeeping (transfer of the margin condition
from $\wstar$ to $\what$) mirrors the local analysis in the companion note
and is deferred to the merged paper.  The formula makes the division of
labor explicit: $Q$ is fixed by the world and bounds what weights can ever
be learned; $b$ is compute, and its exponent is dictated by $\beta$ alone.

\section{What no compute can buy: the diversity floor}
\label{sec:floor}

All results here concern the excess over $R(\wfix)$.  The risk $R(\wfix)$
itself is bounded below by the mass of problems on which the model pool is
jointly hopeless --- e.g.\ no convex combination of the models' answer
distributions puts the gold answer in the lead (on AIME2024, problems where
no policy has gold as plurality).  That floor is a property of model
\emph{diversity}: it is set by the correlation structure of the models'
errors, not by compute, and it caps what any allocation policy can achieve.
A quantitative theory of this floor --- predicting ensemble risk from an
error-dependence structure, an ``effective number of independent models,''
and hence the marginal value of adding a model to the pool --- is the natural
sequel to this work and is deliberately out of scope here.

\section{Empirical program}
\label{sec:empirics}

The theory suggests three experiments, all runnable from saved generations
by extending the bootstrap harness of our resampling study (no new model
calls; semi-synthetic in the same sense as the $N$-floor validation):

\begin{enumerate}
\item \textbf{Uniform-versus-sequential at equal spend.}  Treat the $N=100$
empirical answer distributions as ground truth; run the sequential rule of
Section~\ref{sec:sequential} in the bootstrap and compare its excess-versus-
average-spend curve with uniform allocation at fixed $w$.  Prediction:
uniform decays with slope $-\beta/2$, sequential with slope
$-\beta/(2-\beta)$, on log-log axes (band-limited over the resolvable margin
range, as always on $30$ atoms).
\item \textbf{Spend histogram.}  The sequential rule's per-problem spend
should track $\min\{16L/M_q^2,\nmax\}$: a direct, visualizable signature of
the $m^{-2}$ law on real margins.
\item \textbf{Boundary-atom stress test.}  Run the same comparison at the
raw sign-MILP weights (atoms at $M\approx10^{-4}$): the sequential rule
should be forced to the cap on the atom problems, degrading toward the
uniform rate --- the budget-poison phenomenon of Section~\ref{sec:twosided}
on real data.
\end{enumerate}

Controlled simulations with planted $\beta$ (both channels) then validate
the exponents proper, which no $30$-problem benchmark can identify; the
same simulations exhibit the $\beta=2$ transition (log-excess linear in
$\sqrt b$ at criticality, in $b$ beyond it) and, on finite atomic
populations, the transient-to-exponential crossover at $b\approx\mmin^{-2}$
together with the rate-constant ratio of
Proposition~\ref{prop:hard-margin}.

\appendix

\section{Deferred detail: constants in Corollary~\ref{cor:seq-rate}}
With $\nmax=\lceil b^{2/(2-\beta)}\rceil$, $\delta=1/\nmax$:
$J\asymp\frac{2}{2-\beta}\log b$, $L\le c(\beta,A)\log b$;
the dominant budget term is
$16L\cdot\frac{3C_m}{2-\beta}(\bar m')^{\beta-2}
\asymp(\log b)^{\beta/2}b$ since
$(\bar m')^{\beta-2}=(16L/\nmax)^{(\beta-2)/2}$; the excess band is
$C_m\bar m^\beta=C_m(8L/\nmax)^{\beta/2}
\asymp(\log b)^{\beta/2}b^{-\beta/(2-\beta)}$ and
$\delta=b^{-2/(2-\beta)}$ is of smaller order for $\beta<2$.  Substituting
the reduced budget $b/(C(\log b)^{\beta/2})$ multiplies the excess by
$(\log b)^{\frac{\beta}{2-\beta}\cdot\frac{\beta}{2}\cdot\frac{2-\beta}{\beta}}
=(\log b)^{\beta/2}$-type factors that collapse into the stated
$(\log b/b)^{\beta/(2-\beta)}$ after enlarging the constant.
\end{document}
